\documentclass[answers,12pt,addpoints]{exam} 
\usepackage{import}

\import{C:/Users/prana/OneDrive/Desktop/MathNotes}{style.tex}

% Header
\newcommand{\name}{Pranav Tikkawar}
\newcommand{\course}{Spatial Statistics - STAT 669}
\newcommand{\assignment}{Homework 7.3}
\author{\name}
\title{\course \ - \assignment}

\begin{document}
\maketitle

\newpage
\begin{questions}
\question 8.17

\begin{solution}
The two covariance functions are
\[
K_0(t)=e^{-|t|},\qquad K_1(t)=(1-|t|)_+,
\]
where $(x)_+=\max\{x,0\}$, and $P_j$ denotes the law of a mean-zero
stationary Gaussian process on $[0,L]$ with covariance $K_j$, for
$j=0,1$.  Take $L>1$ and $\delta_n=(L-1)/n$, and define
\[
S_n=\sum_{j=1}^n
\bigl\{Z(\delta_n j)-Z(\delta_n(j-1))\bigr\}
\bigl\{Z(1+\delta_n j)-Z(1+\delta_n(j-1))\bigr\}.
\]

Set $X_{j,n}=Z(\delta_n j)-Z(\delta_n(j-1))$ and
$Y_{j,n}=Z(1+\delta_n j)-Z(1+\delta_n(j-1))$, so
$S_n=\sum_{j=1}^n X_{j,n}Y_{j,n}$.  Since $EZ=0$,
\[
E_j S_n = \sum_{k=1}^n \Cov_j(X_{k,n}, Y_{k,n}).
\]
By stationarity $\Cov_j(X_{k,n},Y_{k,n})$ is the same for every $k$,
and expanding it in terms of $K_j$ gives
\begin{align*}
\Cov_j(X_{k,n},Y_{k,n})
&= E_j\bigl[(Z(\delta_n k)-Z(\delta_n(k-1)))
            (Z(1+\delta_n k)-Z(1+\delta_n(k-1)))\bigr]\\
&= 2K_j(1)-2K_j(1+\delta_n),
\end{align*}
hence $E_j S_n = 2n\{K_j(1)-K_j(1+\delta_n)\}$.

For $K_0(t)=e^{-|t|}$, a direct computation gives
\[
K_0(1+\delta)-K_0(1) = e^{-(1+\delta)}-e^{-1}
= e^{-1}(e^{-\delta}-1)
= -e^{-1}\delta + O(\delta^2),
\]
so $K_0(1)-K_0(1+\delta_n)=e^{-1}\delta_n+O(\delta_n^2)$ and
\[
E_0 S_n = 2n\bigl(e^{-1}\delta_n+O(\delta_n^2)\bigr)
= 2e^{-1}(L-1)+O\!\left(\tfrac{(L-1)^2}{n}\right)
\;\longrightarrow\; 2e^{-1}(L-1).
\]

For $K_1(t)=(1-|t|)_+$, since $\delta_n < 1$ we have $K_1(1)=0$ and
$K_1(1+\delta_n)=0$, so $E_1 S_n = 0$ for all $n$.

Since $2e^{-1}(L-1)>0$, the two sequences have distinct limits.
It remains to show the variances vanish so that $S_n$ concentrates
around each limit.

Both covariance functions satisfy $K_j(0)-K_j(\delta)=|\delta|+O(\delta^2)$
as $\delta\to0$, and since $X_{k,n}$ and $Y_{k,n}$ are both increments
of the same stationary process over an interval of length $\delta_n$,
\[
\sigma_{j,n}^2 := \Var_j(X_{k,n})=\Var_j(Y_{k,n})
= 2\{K_j(0)-K_j(\delta_n)\} = 2\delta_n+O(\delta_n^2).
\]

Write
\[
\Var_j S_n
=\sum_{k=1}^n \Var_j(X_{k,n}Y_{k,n})
+\sum_{k\ne\ell}\Cov_j(X_{k,n}Y_{k,n},X_{\ell,n}Y_{\ell,n}).
\]
Since $Z$ is Gaussian, Isserlis' theorem expresses every fourth-order
joint moment as a sum of products of second-order covariances.
For the diagonal terms: $\Var_j(X_{k,n}Y_{k,n})$ depends only on
$\Var_j(X_{k,n})$, $\Var_j(Y_{k,n})$, and $\Cov_j(X_{k,n},Y_{k,n})$,
each of order $\delta_n$, so $\Var_j(X_{k,n}Y_{k,n})=O(\delta_n^2)$.
For $k\ne\ell$: every term in the Isserlis expansion of
$\Cov_j(X_{k,n}Y_{k,n},X_{\ell,n}Y_{\ell,n})$ contains at least one
cross-index covariance, e.g.\ $\Cov_j(X_{k,n},X_{\ell,n})$.
The lag between $X_{k,n}$ and $X_{\ell,n}$ is $\delta_n|k-\ell|\ge\delta_n>0$;
explicitly,
\[
\Cov_j(X_{k,n},X_{\ell,n})
= K_j\bigl(\delta_n(k-\ell+1)\bigr) - K_j\bigl(\delta_n(k-\ell)\bigr)
- K_j\bigl(\delta_n(k-\ell)\bigr) + K_j\bigl(\delta_n(k-\ell-1)\bigr),
\]
a second difference of $K_j$ at a nonzero lag, which is $O(\delta_n^2)$.
Hence each off-diagonal covariance is $O(\delta_n^3)$, giving
\[
\Var_j S_n
= O(n\delta_n^2)+O(n^2\delta_n^3)
= O\!\left(\frac{(L-1)^2}{n}\right)
 +O\!\left(\frac{(L-1)^3}{n}\right)
\;\longrightarrow\; 0.
\]

By Chebyshev's inequality, $S_n\xrightarrow{P_0}2e^{-1}(L-1)$ and
$S_n\xrightarrow{P_1}0$ in probability.
Setting $a=e^{-1}(L-1)$, which lies strictly between $0$ and
$2e^{-1}(L-1)$, the sets $A_n=\{S_n>a\}$ satisfy
$P_0(A_n)\to1$ and $P_1(A_n)\to0$, so $P_0\perp P_1$.
\end{solution}

\question 8.18

\begin{solution}
The Matérn spectral density in $\mathbb{R}^d$ is
\[
f_j(\omega)
= \frac{\phi_j}{\Bigl(\alpha_j^2
  + \tfrac{1}{2\kappa_j}\,\omega^T A_j \omega\Bigr)^{\kappa_j + d/2}},
\]
with $\phi_j>0$, $\alpha_j>0$, $\kappa_j>0$, and $A_j$ symmetric with
$\operatorname{tr}(A_j)=0$.  We determine when
\[
\int_{|\omega|>R}
\Bigl(1-\frac{f_1(\omega)}{f_0(\omega)}\Bigr)^2\,d\omega < \infty
\]
holds for some $R<\infty$.

As $|\omega|\to\infty$ the $\alpha_j^2$ terms are negligible compared to
the quadratic form, so
\[
f_j(\omega)\sim
\phi_j\Bigl(\tfrac{1}{2\kappa_j}\omega^T A_j\omega\Bigr)^{-(\kappa_j+d/2)}.
\]

Suppose first that at least one parameter other than $\alpha$ differs.

If $A_1\ne A_0$, since $A_1$ and $A_0$ are distinct positive-definite
matrices there exists a direction $\theta_0\in S^{d-1}$ such that
$\theta_0^T A_1\theta_0\ne\theta_0^T A_0\theta_0$.
Along the ray $\omega=r\theta_0$,
\[
\frac{f_1(\omega)}{f_0(\omega)}
\sim
\frac{\phi_1(2\kappa_1)^{\kappa_1+d/2}}{\phi_0(2\kappa_0)^{\kappa_0+d/2}}
\cdot
\frac{(\theta_0^T A_0\theta_0)^{\kappa_0+d/2}}
     {(\theta_0^T A_1\theta_1)^{\kappa_1+d/2}}
\cdot r^{2(\kappa_0-\kappa_1)},
\]
which converges to a limit $C(\theta_0)\ne1$ when $\kappa_0=\kappa_1$,
or diverges/vanishes otherwise.  In either case $(1-f_1/f_0)^2$
does not tend to $0$ along $\omega=r\theta_0$, so
\[
\int_{|\omega|>R}\Bigl(1-\frac{f_1}{f_0}\Bigr)^2\,d\omega
\ge \int_R^\infty\!\int_{B(\theta_0,\varepsilon)}\Bigl(1-\frac{f_1}{f_0}\Bigr)^2
r^{d-1}\,d\sigma(\theta)\,dr = \infty
\]
for any small ball $B(\theta_0,\varepsilon)\subset S^{d-1}$.

If $A_1=A_0=A$ but $\kappa_1\ne\kappa_0$, write $Q(\omega)=\omega^TA\omega$
and $r=|\omega|$.  Then $f_1(\omega)/f_0(\omega)\sim C\,r^{-2(\kappa_1-\kappa_0)}$.
For $\kappa_1>\kappa_0$ the ratio tends to $0$, so $(1-f_1/f_0)^2\to1$
and $\int_{|\omega|>R}(\cdots)\,d\omega\ge c\int_R^\infty r^{d-1}dr=\infty$.
For $\kappa_1<\kappa_0$ the ratio grows like $r^{2(\kappa_0-\kappa_1)}$,
so $(1-f_1/f_0)^2\sim C^2 r^{4(\kappa_0-\kappa_1)}\to\infty$, and the
integral diverges even faster.

If $A_1=A_0$, $\kappa_1=\kappa_0=\kappa$, but $\phi_1\ne\phi_0$, then
$f_1/f_0\to\phi_1/\phi_0\ne1$ uniformly in $\omega$, so
$(1-f_1/f_0)^2\to(1-\phi_1/\phi_0)^2>0$ and (8.16) fails.

So (8.16) requires all parameters except possibly $\alpha$ to be equal.

Now take $\phi_1=\phi_0=\phi$, $\kappa_1=\kappa_0=\kappa$,
$A_1=A_0=A$, but $\alpha_1\ne\alpha_0$.  Then
\[
1-\frac{f_1(\omega)}{f_0(\omega)}
=1-\left(\frac{\alpha_0^2+\tfrac{1}{2\kappa}Q(\omega)}
              {\alpha_1^2+\tfrac{1}{2\kappa}Q(\omega)}\right)^{\kappa+d/2}.
\]
Set $u=\tfrac{1}{2\kappa}Q(\omega)$; for large $|\omega|$,
\[
\frac{\alpha_0^2+u}{\alpha_1^2+u}
=1+\frac{\alpha_0^2-\alpha_1^2}{\alpha_1^2+u}
=1+O(u^{-1}),
\]
and $(1+x)^s=1+sx+O(x^2)$ for small $x$ gives
\[
1-\frac{f_1(\omega)}{f_0(\omega)}
= -(\kappa+d/2)\cdot\frac{\alpha_0^2-\alpha_1^2}{\alpha_1^2+u}+O(u^{-2})
= O(u^{-1})=O\!\bigl(|\omega|^{-2}\bigr).
\]
So $(1-f_1/f_0)^2=O(|\omega|^{-4})$, and in polar coordinates
\[
\int_{|\omega|>R}\Bigl(1-\frac{f_1}{f_0}\Bigr)^2\,d\omega
\;\asymp\;\int_R^\infty r^{d-1}\cdot r^{-4}\,dr
=\int_R^\infty r^{d-5}\,dr,
\]
which converges if and only if $d\le3$.

For $d\le3$, (8.16) holds for two Matérn spectral densities if and only
if all parameters other than $\alpha$ agree.  For $d\ge4$, changing only
$\alpha$ already makes the integral diverge, so (8.16) holds only when
all parameters $(\phi,\alpha,\kappa,A)$ are the same.
\end{solution}

\end{questions}

\end{document}
